% Editable supplement chapter.
\section{Algorithms}
\label{supplement-algorithms}

\subsection{Shared induction strategy}
\label{algorithms}

\label{shared-induction-strategy}
All three learners use separate-and-conquer (sequential covering). They grow a
conjunctive rule greedily, retain it if it covers new observations, remove its
coverage from the working uncovered set, and repeat. The full training set
remains available for evaluating rule quality. Removing an observation from the
uncovered set does not remove it from subsequent quality calculations.

Candidate conditions test nominal equality or numerical thresholds between
adjacent distinct observed values. Candidate generation is local to the
currently covered data. Classification optionally restricts numerical cuts to
class boundaries. Minimum coverage prevents a new rule from simply repeating
previous coverage; {\ttfamily\small max\_growing} limits CR and RR premise lengths separately.
Consequently an ER conjunction may contain more conditions than either parent.

The pseudocodes below describe the methods. They expose the search and acceptance
criteria, rather than serving as line-by-line replacements for the Python
implementation. 
Here {\ttfamily\small Cov(A)} denotes premise coverage in the full training set, {\ttfamily\small U} is the
working uncovered set, and an empty premise initially covers all observations.

\noindent\begin{minipage}{\linewidth}
\textbf{Algorithm 1. Sequential covering (common structure).}


\begin{Verbatim}[fontsize=\scriptsize,frame=single,samepage=true]
INDUCE(D, task, mincov, max_growing):
    R <- empty list
    targets <- class labels for classification; one task for other problems
    for target in targets:
        U <- target-class observations for classification; all D otherwise
        while U is not empty:
            (CR, RR, ER) <- GROW_WITH_EXCEPTION_SEARCH(D, U, target)
            if CR has no conditions: break
            apply task-specific finalisation (see pruning policy below)
            newly_covered <- eligible coverage of CR intersect U
            if newly_covered is empty: break
            append (CR, RR, ER) to R
            U <- U minus newly_covered
    update rule conclusions, coverage statistics, and model metadata
    return R
\end{Verbatim}
\end{minipage}
Classification removes covered positives of the current class from {\ttfamily\small U}.
Regression and survival remove covered observations. A rule without an accepted
exception is still a valid output. Candidate coverage checks also allow final
small remainders in the implementation; {\ttfamily\small mincov} is not a global lower bound
on every reported subgroup size.



\subsection{Classification}
\label{classification-ceris}
The operational classification constraints are
\[
\operatorname{conf}(A\to X)>\theta,\qquad
\operatorname{conf}(B\to X)>\theta,\qquad
\operatorname{conf}(A\land B\to Y)>\theta,\quad Y\ne X.
\]
guarantee in the reported experiments.
The CR predicts a fixed target class \(X\) and grows according to the
configured quality measure. If it is sufficiently precise but still covers
other classes, the learner searches for an RR predicting the same class.
The intersection may predict another class with high precision.

\noindent\begin{minipage}{\linewidth}
\textbf{Algorithm 2. Classification commonsense-rule growth.}


\begin{Verbatim}[fontsize=\scriptsize,frame=single,samepage=true]
GROW_CLASSIFICATION(D, U, X):
    A <- empty premise; RR <- none; ER <- none
    while the growth limit has not been reached:
        c <- highest-quality admissible condition extending A
        if c does not exist: break
        A <- A AND c
        if A covers no observations outside class X: break
        if precision(A -> X) > theta and exception search is enabled:
            (RR, ER) <- SEARCH_CLASSIFICATION_REFERENCE(A, D, X)
            if ER exists: break
    finalise the CR premise using the task's quality history
    return (A -> X, RR, ER)
\end{Verbatim}
\end{minipage}
RR search must simultaneously support \(X\) outside the CR and retain some
of the CR's negatives so that the conjunction can express a class reversal. The current implementation
uses the normalised score


\[
s(B)=q(B\to X)
   +\frac{|C_B\cap U_X|}{|U_X|}
   +\frac{|C_B\cap N_A|}{|N_A|},
\]
where \(N_A=\{i\in C_A:y_i\ne X\}\) and \(U_X\) is the reference
search's set of uncovered positives outside the CR. A ratio with an empty
denominator contributes zero. 

\noindent\begin{minipage}{\linewidth}
\textbf{Algorithm 3. Classification reference and exception search.}


\begin{Verbatim}[fontsize=\scriptsize,frame=single,samepage=true]
SEARCH_CLASSIFICATION_REFERENCE(A, D, X):
    B <- empty premise
    U_X <- class-X observations outside A
    N_A <- non-X observations covered by A
    repeat until growth limit or no eligible refinement:
        evaluate candidate extensions B' using reference score s(B')
        require new positive coverage and precision(B' -> X) > theta
        test the intersection A AND B' for an alternative class Y
        prefer valid exception candidates using the selector's score rules
        extend B with the selected condition
        if the selected intersection predicts Y != X with precision > theta:
            return (B -> X, A AND B -> Y)
    return (none, none)
\end{Verbatim}
\end{minipage}
The alternative label is the modal class among the intersection's non-X
observations, and its precision is calculated on the entire intersection.
The selector can favour an admissible exception candidate before it improves
the previous score; it is not simply a global maximisation of \(s\).
An empty intersection of negatives cannot produce an attached exception.


\subsection{Regression}
\label{algorithm-regression}
To construct rule positive/negative counts,
the learner calculates the covered target mean and standard deviation. Targets
inside this local mean-plus/minus-standard-deviation band are treated as
positive, and the same band defines dataset-wide positive totals. This is a
quality calculation for numerical targets, rather than conversion to fixed
class labels.

For the statistical screen, let $E=C_A\cap C_B$, $G_A=C_A\setminus E$,
and $G_B=C_B\setminus E$. The three groups are disjoint. The current code
accepts a candidate only when the two-sided Mann--Whitney comparisons satisfy
\[
p(G_A,G_B)>0.05,\qquad p(E,G_A)\leq0.05,\qquad p(E,G_B)\leq0.05.
\]


\noindent\begin{minipage}{\linewidth}
\textbf{Algorithm 4. Growing a regression commonsense rule.}


\begin{Verbatim}[fontsize=\scriptsize,frame=single,samepage=true]
GROW_REGRESSION(D, U):
    A <- empty premise
    while the growth limit has not been reached:
        c <- highest-C2 admissible condition extending A
        if c does not exist: break
        A <- A AND c
        update covered-target mean and variability band
        if exception search is enabled:
            (RR, ER) <- SEARCH_REGRESSION_REFERENCE(A, D)
            if ER exists: return (A, RR, ER)
    return (A, none, none)
\end{Verbatim}
\end{minipage}
\noindent\begin{minipage}{\linewidth}
\textbf{Algorithm 5. Regression reference search.}


\begin{Verbatim}[fontsize=\scriptsize,frame=single,samepage=true]
SEARCH_REGRESSION_REFERENCE(A, D):
    B <- empty premise; best <- none; best_overlap <- 0
    repeat until growth limit:
        selected <- none
        for each unused elementary condition c extending B:
            B' <- B AND c; E' <- A AND B'
            require sufficient new reference coverage outside A
            require compatible parent means and exceptional intersection mean
            require the three disjoint-group Mann-Whitney criteria
            if |Cov(E')| > best_overlap:
                selected <- c; best <- (B', E')
                best_overlap <- |Cov(E')|
        if selected is none: break
        B <- B AND selected
    return best, or (none, none)
\end{Verbatim}
\end{minipage}

\subsection{Survival analysis}
\label{algorithm-survival-analysis}
Survival CR conditions maximise log-rank separation between covered and
uncovered observations. RR candidates are feasible if the CR and RR comparison
does not reject at \(0.05\), while both comparisons involving their
intersection reject at that level. Among feasible candidates, the reference
selector favours a larger CR--RR p-value. This ranking heuristic is not an
equivalence test.

\noindent\begin{minipage}{\linewidth}
\textbf{Algorithm 6. Survival CR and reference search.}


\begin{Verbatim}[fontsize=\scriptsize,frame=single,samepage=true]
GROW_SURVIVAL(D, U):
    A <- empty premise
    repeat until growth limit:
        c <- admissible condition with highest covered-vs-uncovered log-rank
        if c does not exist: break
        A <- A AND c
        if exception search is disabled: continue
        B <- empty premise; best_score <- 0
        repeat until reference growth limit:
            for candidate B' = B AND c_ref:
                E' <- A AND B'
                fit Kaplan-Meier curves on Cov(A), Cov(B'), Cov(E')
                calculate p_AB, p_AE, p_BE by log-rank comparisons
                require nonempty E' and sufficient new coverage outside A
                require p_AB > 0.05 and p_AE <= 0.05 and p_BE <= 0.05
                retain candidate if p_AB improves best_score
            if no candidate retained: break
            B <- retained candidate; update best_score
        finalise B using the implementation's quality history
        recheck intersection separation before attaching RR and ER
        if accepted: break
    finalise A and fit its Kaplan-Meier conclusion
    return (A, attached RR or none, attached ER or none)
\end{Verbatim}
\end{minipage}
Prediction uses the KM estimator of an applicable exception in place of its
parent rule's estimator, then averages the estimators contributed by covering
rules. The source of each prediction remains
inspectable through the rule structure.




